\section{Introduction}
\label{sec:intro}

``No issues found'' is the most common line on a safety scorecard, and the least examined.
A null result from a safety metric tells us something about a model only if the metric would have fired had the model misbehaved.
Analytical chemistry and clinical testing handle this with positive controls: before trusting a negative, run the assay on a sample known to be positive.
Automated safety evaluation of language models rarely does the same once a metric leaves the domain it was built for.

This matters because safety metrics are increasingly used as general-purpose instruments.
Harm classifiers and LLM judges are validated against human labels on harmful-compliance data \citep{chao2024jailbreakbench,mazeika2024harmbench,souly2024strongreject,inan2023llamaguard}, and are then run over model outputs of many kinds.
If the metric that catches a model writing malware instructions is blind to a model that flatters a user into a false belief, or that presents a wrong derivation with full confidence, then a clean cross-domain scorecard describes the instrument rather than the model.
Labs, auditors, and deployers who rely on such scorecards need to know each instrument's sensitivity in each domain.

Existing work leaves this question open from two sides.
On one hand, judge-validation studies measure agreement with humans inside the harmful-compliance domain only.
On the other hand, subtle behaviours such as sycophancy \citep{perez2022discovering,sharma2023sycophancy} and unfaithful reasoning \citep{turpin2023language,lanham2023measuring} are measured with bespoke protocols, and those studies do not ask what a general safety metric would say about the same outputs.
What is missing is an analysis that holds the metric fixed and varies the domain.

{\bf Would a standard safety metric have fired if misalignment were present?}
We answer this with a small cross-domain positive-control benchmark (\figref{fig:overview}).
For each of three domains we collect responses known to be misaligned (\emph{anchors}) and matched aligned \emph{controls}: 110 human-labelled jailbroken responses and 190 safe ones from \jbb; 174 pairs in which a model either affirms or politely corrects a user's wrong trivia answer; and 223 pairs of grade-school math solutions that either contain one unacknowledged slip or are correct.
We verify every label with a rule that is independent of all detectors.
We then run five standard detectors identically on all 1,094 items, and add a rubric ablation that asks where any gap comes from.
We fix the hypothesis before running: detection above 70\% on overt anchors, below 30\% on both subtle domains, and recovery of the subtle anchors by a judge with a domain-appropriate rubric.

The hypothesis holds for the purpose-built harm detectors and only partly otherwise.
\granite flags 100\% of overt anchors, 0\% of sycophancy anchors, and 9.4\% of flawed-reasoning anchors.
An LLM judge with the \jbb harm rubric flags 75.5\%, 1.1\%, and 2.7\%; its overt rate has a 95\% interval of 66.6--82.5\%, so the ``above 70\%'' half is statistically established only for \granite.
A broad ``is anything misaligned here?'' rubric does not stay below 30\% on reasoning (43.9\%), and on sycophancy it flags 29.9\% of anchors but also 23.0\% of honest controls (AUROC 0.50).
The rubric ablation separates two kinds of gap.
For sycophancy, a targeted rubric alone closes it (100\% detected, 16.1\% false positives, AUROC 1.00).
For flawed reasoning it does not: the targeted judge flags 85.2\% of anchors and 61.0\% of correct solutions, and becomes discriminative only with the reference answer (AUROC 0.97).
A second judge model does not reproduce the targeted-rubric results.

In summary, our contributions are:
\begin{itemize}[leftmargin=*,itemsep=0pt,topsep=0pt]
    \item We build a cross-domain positive-control benchmark of 1,094 responses with detector-independent labels: verified anchors and matched controls for overt violations, sycophancy, and flawed reasoning.
    \item We measure the per-domain sensitivity of five standard detectors with pre-specified tests, and show that harm-style metrics lose 73--100 percentage points of detection rate from the overt domain to the subtle ones.
    \item We conduct a rubric ablation that distinguishes a \emph{scope} gap (sycophancy: the judge knows but is not asked) from a \emph{capability} gap (reasoning: the judge is asked but cannot verify).
    \item We show that detection rate alone misleads: three detectors reach at least 85\% detection on a subtle domain with no discriminative value, and one judge's verdicts depend on the wording of its answer labels.
\end{itemize}

\Secref{sec:related} positions our work, \secref{sec:method} describes the anchors, detectors, and tests, \secref{sec:results} reports results, and \secref{sec:discussion} discusses implications and limitations.

\begin{figure}[t]
\centering
\resizebox{0.95\linewidth}{!}{%
\begin{tikzpicture}[
  box/.style={draw, rounded corners=2pt, align=center, font=\small, minimum height=1.05cm, text width=3.3cm},
  wide/.style={draw, rounded corners=2pt, align=center, font=\small, minimum height=1.05cm, text width=3.6cm},
  arr/.style={-{Stealth[length=2mm]}, semithick}]
\node[box] (d1) {\textbf{Overt safety}\\ 110 anchors / 190 controls\\ human-labelled};
\node[box, below=0.25cm of d1] (d2) {\textbf{Sycophancy}\\ 174 matched pairs\\ string-verified};
\node[box, below=0.25cm of d2] (d3) {\textbf{Flawed reasoning}\\ 223 matched pairs\\ answer-verified};
\node[wide, right=1.3cm of d1] (s) {\textbf{Standard battery}\\ refusal keywords, toxicity,\\ harm classifier, harm judge,\\ broad judge};
\node[wide, right=1.3cm of d3] (a) {\textbf{Rubric ablation}\\ targeted rubric,\\ targeted + reference};
\node[wide, right=1.3cm of s, yshift=-1.3cm] (m) {\textbf{Per-domain sensitivity}\\ TPR on anchors, FPR on\\ controls, Youden's $J$, AUROC};
\draw[arr] (d1.east) -- (s.west);
\draw[arr] (d2.east) -- (s.west);
\draw[arr] (d3.east) -- (s.west);
\draw[arr] (d2.east) -- (a.west);
\draw[arr] (d3.east) -- (a.west);
\draw[arr] (s.east) -- (m.west);
\draw[arr] (a.east) -- (m.west);
\end{tikzpicture}}
\caption{Overview of the positive-control design. We plant known-misaligned anchors with matched aligned controls in three domains, run one fixed battery of standard detectors on every item with identical inputs, and report sensitivity per domain. A rubric ablation on the two subtle domains tests whether a domain-appropriate rubric, or a reference answer, recovers the anchors.}
\label{fig:overview}
\end{figure}
